\documentclass[11pt]{article}
\PassOptionsToPackage{numbers,sort&compress}{natbib}
\usepackage[final]{acl}
\usepackage{times}
\usepackage{latexsym}
\usepackage[T1]{fontenc}
\usepackage[utf8]{inputenc}
\usepackage{microtype}
\usepackage{amsmath,amssymb}
\usepackage{booktabs,array}
\usepackage[table]{xcolor}
\usepackage{tikz}
\usetikzlibrary{arrows.meta,positioning}
\usepackage{float}
\usepackage{graphicx}

\definecolor{pipeblue}{RGB}{30,100,180}
\definecolor{pipeteal}{RGB}{0,140,130}
\definecolor{pipeorange}{RGB}{210,100,20}
\definecolor{pipegreen}{RGB}{30,140,60}
\definecolor{pipegray}{RGB}{80,80,90}
\definecolor{headerblue}{RGB}{220,235,255}
\definecolor{rowtint}{RGB}{245,250,255}

\setlength\titlebox{7\baselineskip}

\title{PrimeLine@DravidianLangTech 2026: Abusive Tamil Comment Detection\\Using MuRIL}

\author{Rithikaa~V \quad Sanjay~Krishnan~K \quad Nithya~Varshini~C~N~R \\
  \textit{Guided by: S.~Sumathi} \\
  Department of Information Technology, St.~Joseph\'s College of Engineering \\
  \url{[GitHub link to be added]}
}

\begin{document}
\maketitle

\begin{abstract}
Detecting abusive language in Tamil social media is a genuinely difficult problem.
The language is morphologically rich, speakers routinely mix Tamil with English,
and informal romanised Tamil is common enough to confuse models trained primarily
on formal text. This work presents a system for binary classification of Tamil
comments into Abusive and Non-Abusive categories, submitted to the
DravidianLangTech@ACL 2026 shared task \citep{abusive-tamil-women-dravidianlangtech-acl-2026}.
MuRIL \citep{khanuja2021}, a BERT-based encoder pre-trained on 17 Indian languages
and their transliterated equivalents, is fine-tuned, and it is shown that this
Indian-language-specific pre-training provides a meaningful advantage over generic
multilingual baselines. The system achieves a macro-averaged F1 of 0.83 on the
validation set, compared to 0.79 for XLM-RoBERTa and 0.77 for mBERT under identical
training conditions, establishing a strong transformer-based baseline for abusive
language detection in code-mixed Tamil.
\end{abstract}

\section{Introduction}

Tamil is one of the oldest classical languages still in everyday use, with more than
80 million speakers worldwide. Social media has given Tamil speakers an enormous
platform for public expression, but it has also made it easier for abusive and harmful
content to spread. Detecting such content automatically is a growing research priority,
yet most existing systems are designed for high-resource languages like English and do
not transfer well to Tamil.

The challenge goes beyond data scarcity. Tamil is an agglutinative language, where a
single word can encode morphological information that would require several words in
English. On social media, this complexity is compounded by heavy code-mixing: users
blend native Tamil script, romanised Tamil, and English freely within a single comment,
with no consistent orthographic convention. Lexicon-based systems and classical machine
learning approaches struggle in this setting because they rely on vocabulary stability
that code-mixed text simply does not have.

The DravidianLangTech@ACL 2026 shared task on abusive Tamil comment detection
\citep{abusive-tamil-women-dravidianlangtech-acl-2026} provides a structured benchmark
for this problem, framing it as binary classification: given a Tamil social media
comment, decide whether it is Abusive or Non-Abusive. The approach centers on
fine-tuning MuRIL \citep{khanuja2021}, a transformer encoder pre-trained specifically
on 17 Indian languages including Tamil, using both native script and transliterated
corpora. MuRIL\textquotesingle{}s dual-script coverage makes it uniquely well-suited
to the romanised Tamil that is pervasive in online communication.

\section{Related Work}

Early systems for abusive language detection relied on keyword lists and hand-crafted
features combined with traditional classifiers such as Support Vector Machines and
Naive Bayes \citep{nobata2016}. While effective for formal single-language text, these
methods break down when vocabulary is informal or code-mixed.

The introduction of BERT \citep{devlin2019} changed the landscape considerably.
Pre-trained transformer models produced contextual representations that generalised
far better across domains, and their multilingual variants extended these benefits to
non-English languages. \newcite{ranasinghe2021} demonstrated that fine-tuned
multilingual transformers consistently outperform traditional methods on offensive
language tasks across multiple languages. \newcite{vaswani2017} introduced the
Transformer architecture that underpins all of these models, establishing multi-head
self-attention as the mechanism for capturing long-range token dependencies.
\newcite{liu2019} later showed that more robust pre-training strategies yield better
downstream performance, a finding extended to the multilingual setting by
\newcite{conneau2020} through XLM-RoBERTa.

For Dravidian languages, the DravidianLangTech and LT-EDI shared tasks have been the
primary source of benchmarks. \newcite{chakravarthi2021offensive} reported results on
offensive language identification in Tamil, Malayalam, and Kannada, where fine-tuned
mBERT and XLM-RoBERTa were the dominant approaches. \newcite{sivasai2021} explored
offensive language identification across Dravidian languages. Tamil-specific work has
explored ensemble strategies and model combinations: \newcite{anbukkarasi2023} applied
deep learning for hate speech detection in code-mixed Tamil, \newcite{subramanian2022}
used adapters and cross-domain transfer for offensive language in Tamil YouTube
comments, and \newcite{krishna2024} addressed class imbalance in Tamil code-mixed
abusive comment detection. \newcite{nishanth2025} applied RoBERTa and XGBoost for
abusive Tamil and Malayalam text detection at DravidianLangTech 2025. However, MuRIL
has not been systematically evaluated on abusive language detection in Tamil ---
which is the gap this work addresses.

MuRIL \citep{khanuja2021} was developed by Google Research and pre-trained on 17
Indian languages using Wikipedia and Common Crawl data in both native and transliterated
script. This dual-script design is a key advantage over XLM-RoBERTa and mBERT, neither
of which includes transliterated training data. For Tamil social media, where romanised
text is widespread, MuRIL\textquotesingle{}s pre-training distribution is a much closer
match to the target domain.

\section{Dataset and Preprocessing}

\subsection{Dataset}
The dataset was provided by the shared task organisers and consists of Tamil social
media comments labelled as Abusive or Non-Abusive. The text reflects real-world Tamil
online writing, freely mixing native Tamil script, romanised Tamil, and English within
individual comments. This code-mixed, multi-script nature makes automated classification
considerably harder than a standard monolingual setting.

The full annotated training set contains 3,652 examples. A stratified 90/10 split is
applied to preserve label distribution, yielding 3,286 training examples and 366
validation examples. The test set comprises 912 unlabelled comments for blind evaluation,
which is standard practice in shared task settings. The training data exhibits moderate
class imbalance: approximately 64\% Non-Abusive and 36\% Abusive. Table~\ref{tab:dist}
shows the full class distribution. Comment lengths vary from brief single-word
utterances to multi-sentence posts, with a mean token length of approximately 18 tokens
after Unicode normalisation. The minority Abusive class is the harder of the two to
detect reliably, especially when abuse is expressed through sarcasm or indirect phrasing
rather than explicit offensive vocabulary.

\begin{table}[H]
\centering\small
\setlength{\tabcolsep}{6pt}\renewcommand{\arraystretch}{1.1}
\rowcolors{2}{rowtint}{white}
\begin{tabular}{lcc}
\toprule
\rowcolor{headerblue}\textbf{Label} & \textbf{Count} & \textbf{\%} \\
\midrule
Non-Abusive & 2338 & 64\% \\
Abusive     & 1314 & 36\% \\
\midrule
\textbf{Total} & \textbf{3652} & \textbf{100\%} \\
\bottomrule
\end{tabular}
\caption{Class distribution for the Tamil abusive comment detection training data.}
\label{tab:dist}
\end{table}

\subsection{Preprocessing}
Text samples are tokenised using the MuRIL WordPiece tokeniser, which handles both
native Tamil script and romanised equivalents natively. A maximum sequence length of
128 tokens is used, sufficient for the short social media comments in this dataset.
Sequences exceeding this limit are truncated, and shorter sequences are padded within
each batch. Unicode normalisation is applied to handle encoding inconsistencies common
in crawled social media text.

Transliteration, language identification, and stopword removal are deliberately skipped.
MuRIL\textquotesingle{}s dual-script pre-training already accounts for the mix of
native and romanised Tamil, and additional normalisation risks erasing code-switching
cues that carry genuine semantic meaning.

\section{Methodology}

\subsection{System Pipeline}
The system follows a straightforward fine-tuning pipeline, illustrated in
Figure~\ref{fig:pipeline}. Each comment is passed to the MuRIL tokeniser, encoded
as token IDs with a prepended \texttt{[CLS]} token, and fed into the MuRIL transformer
encoder. The hidden state at the \texttt{[CLS]} position is extracted and passed through
a linear classification head that maps it to two logits. These are normalised via
softmax, and the class with the highest probability is taken as the prediction.

\begin{figure}[H]
\centering
\begin{tikzpicture}[node distance=3mm,
  box/.style={rectangle,rounded corners=2pt,minimum width=\columnwidth-2mm,
    minimum height=7mm,align=center,font=\small\bfseries,draw,thick},
  arr/.style={-Stealth,thick,color=pipegray}]
\node[box,fill=pipeorange!25,draw=pipeorange!80!black]
  (inp) {Raw Text Input (Code-Mixed Tamil)};
\node[box,fill=pipeteal!20,draw=pipeteal!80!black,below=of inp]
  (tok) {MuRIL Tokeniser (WordPiece, max 128)};
\node[box,fill=pipeblue!20,draw=pipeblue!80!black,below=of tok]
  (enc) {MuRIL Encoder (12 layers, 768 dims)};
\node[box,fill=pipeblue!10,draw=pipeblue!60!black,below=of enc]
  (cls) {\texttt{[CLS]} Token Representation};
\node[box,fill=pipegreen!20,draw=pipegreen!80!black,below=of cls]
  (hd) {Linear Head: 2 classes};
\node[box,fill=pipeorange!30,draw=pipeorange!80!black,below=of hd]
  (out) {Predicted Label (Abusive / Non-Abusive)};
\draw[arr](inp)--(tok);\draw[arr](tok)--(enc);\draw[arr](enc)--(cls);
\draw[arr](cls)--(hd);\draw[arr](hd)--(out);
\end{tikzpicture}
\caption{End-to-end system pipeline for abusive comment classification in Tamil.}
\label{fig:pipeline}
\end{figure}

\subsection{Model: MuRIL}
The core model is MuRIL (Multilingual Representations for Indian Languages;
\citealt{khanuja2021}), a BERT-based encoder released by Google Research, pre-trained
on 17 Indian languages in both native and transliterated script. Its 12-layer
transformer architecture produces 768-dimensional contextualised token representations.
A single linear classification head is placed over the \texttt{[CLS]} token
representation, mapping the 768-dimensional vector to two output logits. The entire
model is fine-tuned end-to-end.

\subsection{Model Formulation}
At each transformer layer, MuRIL computes scaled dot-product attention over queries
$Q$, keys $K$, and values $V$:
\begin{equation}
\text{Attention}(Q,K,V) = \text{softmax}\!\left(\frac{QK^{\top}}{\sqrt{d_k}}\right)V
\end{equation}
where $d_k$ is the key vector dimension. This allows the model to capture dependencies
between tokens across Tamil script, romanised Tamil, and English simultaneously within
a single comment.

The \texttt{[CLS]} representation from the final encoder layer is passed through the
linear head, and the resulting logits are normalised via softmax to produce class
probabilities for $c \in \{\text{Abusive},\,\text{Non-Abusive}\}$:
\begin{equation}
P(y = c \mid \mathbf{x}) = \frac{\exp(z_c)}{\sum_{j}\exp(z_j)}
\end{equation}

Training minimises binary cross-entropy loss over all training samples:
\begin{equation}
\mathcal{L} = -\frac{1}{N}\sum_{i=1}^{N}\Bigl[
  y_i \log P(y{=}1 \mid \mathbf{x}_i)
  + (1-y_i)\log P(y{=}0 \mid \mathbf{x}_i)\Bigr]
\end{equation}
where $N$ is the number of training examples and $y_i \in \{0,1\}$ is the binary
ground truth label.

\subsection{Training Setup}
All models are fine-tuned using the Hugging Face Transformers Trainer API
\citep{wolf2020} with a learning rate of $2 \times 10^{-5}$, batch size of 16,
weight decay of 0.01, and 4 training epochs. The AdamW optimiser is used throughout.
Macro-averaged F1 on the validation set serves as the primary evaluation metric,
since it weights all classes equally and is sensitive to performance on the minority
Abusive class. All experiments run on a CUDA-enabled GPU with total training time of
approximately 25 minutes.

MuRIL was evaluated at 2, 3, 4, and 5 training epochs on the validation set.
Performance improved from 0.79 at 2 epochs to 0.83 at 4 epochs, and stabilised
thereafter with no further gain at 5 epochs. The 4-epoch configuration was therefore
used for all reported experiments.

\section{Results and Discussion}

\subsection{Official Results}
Table~\ref{tab:results} reports macro-averaged precision, recall, and F1 on the
held-out validation set for MuRIL and two multilingual baselines --- mBERT
\citep{devlin2019} and XLM-R \citep{conneau2020} --- both fine-tuned under identical
settings.

\begin{table}[H]
\centering\small
\setlength{\tabcolsep}{5pt}\renewcommand{\arraystretch}{1.1}
\rowcolors{2}{rowtint}{white}
\begin{tabular}{lccc}
\toprule
\rowcolor{headerblue}\textbf{Model} & \textbf{P} & \textbf{R} & \textbf{F1} \\
\midrule
mBERT              & 0.78 & 0.76 & 0.77 \\
XLM-R              & 0.80 & 0.79 & 0.79 \\
\textbf{MuRIL (proposed)} & \textbf{0.84} & \textbf{0.83} & \textbf{0.83} \\
\bottomrule
\end{tabular}
\caption{Macro-averaged results on the validation set (P\,=\,Precision, R\,=\,Recall, F1\,=\,Macro F1).}
\label{tab:results}
\end{table}

MuRIL achieves the best performance across all metrics, with a macro-F1 of 0.83
compared to 0.79 for XLM-R and 0.77 for mBERT. The improvement is most pronounced
for the Abusive class, where MuRIL\textquotesingle{}s Indian-language-specific
pre-training enables better representation of Tamil morphology and colloquial script
variants. The gap between MuRIL and XLM-R is especially telling: XLM-R is a strong
multilingual model, but it does not include transliterated training data. For Tamil
social media, that gap in pre-training coverage translates directly into lower recall
on abusive content.

Table~\ref{tab:perclass} shows the per-class breakdown for MuRIL, giving a clearer
view of where the model performs well and where it still falls short.

\begin{table}[H]
\centering\small
\setlength{\tabcolsep}{6pt}\renewcommand{\arraystretch}{1.1}
\rowcolors{2}{rowtint}{white}
\begin{tabular}{lcc}
\toprule
\rowcolor{headerblue}\textbf{Class} & \textbf{Precision} & \textbf{Recall} \\
\midrule
Non-Abusive & 0.87 & 0.91 \\
Abusive     & 0.81 & 0.75 \\
\bottomrule
\end{tabular}
\caption{Per-class precision and recall for the proposed MuRIL system.}
\label{tab:perclass}
\end{table}

The Non-Abusive class is handled with higher precision and recall, as expected given
the class imbalance. A recall of 0.75 on Abusive comments means roughly one in four
abusive examples is missed --- improving this is the clearest direction for future work.

\subsection{Error Analysis}
Qualitative examination of the misclassified examples reveals two dominant patterns.
The first is implicit abuse: comments that convey abusive intent through sarcasm,
euphemism, or cultural reference without any overtly offensive vocabulary. These are
consistently misclassified as Non-Abusive by all three models, since the meaning only
becomes clear with contextual knowledge the models do not possess.

The second pattern involves dense code-mixing: utterances that switch between Tamil
script, romanised Tamil, and English multiple times within a single short comment.
When all three layers are simultaneously active, even MuRIL occasionally fails to
build a coherent contextual representation.

False negatives (Abusive predicted as Non-Abusive) outnumber false positives across
all models, reflecting the majority-class bias introduced by the training imbalance.
Future work should explore loss reweighting or focal loss to counteract this tendency,
as well as back-translation from related Dravidian languages to augment the Abusive
training set.

\section{Conclusion}
This work presents a system for Tamil abusive comment detection based on fine-tuning
MuRIL for binary sequence classification. By leveraging Indian-language-specific
pre-training on both native and transliterated scripts, the model achieves a macro-F1
of 0.83 on the validation set, outperforming XLM-R and mBERT by clear margins under
identical conditions. The results confirm that pre-training data composition matters:
a model exposed to transliterated Tamil is better equipped to handle the romanised
code-mixing that defines Tamil social media. Future directions include adversarial
training for implicit abuse, ensemble methods combining transformer representations
with lexicon-based features, and extension to multi-class offensive language categories.

\section*{Limitations}
The system is trained and evaluated on the provided shared task dataset alone, and
generalisation to other Tamil platforms or time periods cannot be guaranteed. The
maximum sequence length of 128 tokens may truncate longer comments, losing relevant
context. Computational constraints limited experiments to MuRIL-base. Class imbalance
was not addressed through resampling or weighted loss, which likely suppressed recall
on the Abusive class.

\section*{Acknowledgments}
The organizers of the DravidianLangTech@ACL~2026 shared task
\citep{abusive-tamil-women-dravidianlangtech-acl-2026} are thanked for providing
the annotated dataset and evaluation infrastructure. Sincere gratitude is expressed
to the project guide, S.~Sumathi, Department of Information Technology,
St.~Joseph\textquotesingle{}s College of Engineering, for constant guidance and
encouragement throughout this work. The developers of MuRIL and the Hugging Face
Transformers library are also thanked for making their tools openly available.

\bibliography{tamil_custom}

\end{document}